%% Experiences section.

\section{\faBriefcase$\;$\textbf{Expériences professionnelles}}
\cventry{Février 2018 $\rightharpoonup$\\Maintenant}{\href{https://www.wemanity.com}{Wemanity}}{Agile Consultant Developer R\&D}{Paris (10ème)}{}{%
  \ifthenelse{\equal{\displayImageExperiences}{yes}}{
    \vspace{-0.2cm}
    \begin{wrapfigure}{R}{0.2\textwidth}%
      \vspace{-1.1cm}
      \centering%
      \href{https://www.dhatim.com}{\includegraphics[width=0.18\textwidth]{./img/wemanity}}%
    \end{wrapfigure}%
  }{}
  Développeur JavaScript fullstack au sein du Lab Innovation \textit{Innov8} de la \href{https://www.societegenerale.fr/}{Société Générale}. En charge de développer des PoC innovants autour de la 3D/VR pour de la modélisation de données complexes. Etude et développement d'applications autour des robots \href{https://www.softbankrobotics.com/emea/en/robots/pepper}{\textit{Pepper}} de \href{https://www.softbankrobotics.com/emea/en}{Softbank Robotics} pour des évènements et cas d'usage internes.\\Stack technique: ES6, ReactJS, Three.js, Loopback, Python, Docker, OpenShift.}

\vspace{\cvExperienceSpace}

\cventry{Octobre 2016 $\rightharpoonup$\\Février 2018}{\href{https://www.dhatim.com}{Dhatim}}{Développeur Java R\&D}{Paris (12ème)}{}{%
  \ifthenelse{\equal{\displayImageExperiences}{yes}}{
    \vspace{-0.2cm}
    \begin{wrapfigure}{R}{0.2\textwidth}%
      \vspace{-0.3cm}
      \centering%
      \href{https://www.dhatim.com}{\includegraphics[width=0.18\textwidth]{./img/dhatim}}%
    \end{wrapfigure}%
  }{}
  Développeur backend attaché à l'équipe de \href{https://pay.conciliator.ai/}{\textit{Conciliator For Pay}}, solution en SaaS de contrôle de DSN (Déclaration Sociale Nominative) pour les entreprises. Équipe composée de 5 personnes: 2 front, 2 back, 1 PO en méthode Scrum.\\Stack technique: Java 8 Core, Dropwizard, PostgreSQL, Docker, Git, Jenkins.}

\vspace{\cvExperienceSpace}

\cventry{Mars 2015 $\rightharpoonup$\\Octobre 2016}{\href{https://www.visian.fr}{Visian}}{Lead Developper \& Responsable Fablab}{Nanterre (92)}{}{%
  \ifthenelse{\equal{\displayImageExperiences}{yes}}{
    \vspace{-0.2cm}
    \begin{wrapfigure}{R}{0.2\textwidth}%
      \vspace{-0.7cm}
      \centering%
      \href{https://www.visian.fr}{\includegraphics[width=0.18\textwidth]{./img/visian}}%
    \end{wrapfigure}%
  }{}
  Lead Dev \& Lead Tech chez Visian, la BU Innovation autour des objets connectés de \href{https://www.neurones-it.com}{Neurones IT}. Responsable de l'équipe des développeurs (hard \& soft). Réalisation et développement de PoC en méthodes agiles mêlant IoT (beacons, wearables, Raspberry, tags NFC, capteurs...) et cloud pour la remontée de données (Microsoft Azure).}

\vspace{\cvExperienceSpace}

\cventry{Septembre 2014 $\rightharpoonup$\\Mars 2015}{\href{https://www.groupe-sii.com/fr/}{Groupe SII}}{Assistant Responsable Innovation}{Paris (75)}{}{%
  \ifthenelse{\equal{\displayImageExperiences}{yes}}{
    \vspace{-0.2cm}
    \begin{wrapfigure}{R}{0.2\textwidth}%
      \vspace{-0.8cm}
      \centering%
      \href{https://www.groupe-sii.com/fr/}{\includegraphics[width=0.18\textwidth]{./img/sii}}%
    \end{wrapfigure}%
  }{}
  Étude et exploration des drones Open Source tels que le projet \href{https://www.bitcraze.se/}{Crazyflie}. Réalisation d’un prototype d'objet connecté à partir d’un \href{https://raspberrypicomputer.com/}{Raspberry Pi} et de divers capteurs (température, pression, GPS). Développement avec le SDK Java 8 Embedded ME for IoT, et remontée des données.}

\vspace{\cvExperienceSpace}

\cventry{Avril 2013 $\rightharpoonup$\\Août 2013}{\href{https://www.thalesgroup.fr/}{Thales Services}}{Stagiaire Administrateur Système}{Ramonville (31)}{}{%
  \ifthenelse{\equal{\displayImageExperiences}{yes}}{
    \vspace{-0.2cm}
    \begin{wrapfigure}{R}{0.2\textwidth}%
      \vspace{-0.3cm}
      \centering%
      \href{https://www.thalesgroup.fr/}{\includegraphics[width=0.18\textwidth]{./img/thales}}%
    \end{wrapfigure}%
  }{}
  Responsable de l'étude, de l'installation et du déploiement d'\href{https://www.openstack.org}{OpenStack} dans un environnement de tests. Rédaction du dossier d’architecture matérielle, réalisation des scripts de déploiement automatique, tests et validations.}

\vspace{\cvExperienceSpace}

\cventry{Octobre 2012 $\rightharpoonup$\\Avril 2013}{\href{https://www.novacom-services.com/}{Novacom Services}}{Stagiaire Développeur}{Ramonville (31)}{}{%
  \ifthenelse{\equal{\displayImageExperiences}{yes}}{
    \vspace{-0.2cm}
    \begin{wrapfigure}{R}{0.2\textwidth}%
      \vspace{-0.8cm}
      \centering%
      \href{https://www.novacom-services.com/}{\includegraphics[width=0.18\textwidth]{./img/novacom}}%
    \end{wrapfigure}%
  }{}
  Développement d'un jeu de données et d'un ensemble de tests unitaires et fonctionnels sur la couche service des applications, permettant l’automatisation des validations logicielles et ainsi garantir la non régression pendant les développements.}

\vspace{\cvSectionSpace}
